% ============================================================
%  从Transformer数学结构到Agent工程约束：一个收敛理论框架
%  arXiv 完整稿 v2
%  整合闭合分析、硬伤修复、形式化推导、全面文献对话与并行研究验证
% ============================================================

\documentclass[11pt,a4paper]{ctexart}

% ---- 页面 ----
\usepackage[left=3cm,right=3cm,top=2.5cm,bottom=2.5cm]{geometry}

% ---- 数学 ----
\usepackage{amsmath,amssymb,amsthm,mathtools}

% ---- 图表 ----
\usepackage{graphicx,booktabs,array,caption,subcaption}
\usepackage{tikz}
\usetikzlibrary{arrows,shapes,positioning,fit,backgrounds,chains}

% ---- 链接 ----
\usepackage[colorlinks=true,linkcolor=blue,citecolor=blue,urlcolor=blue]{hyperref}
\usepackage{url}

% ---- 定理 ----
\newtheorem{theorem}{定理}[section]
\newtheorem{proposition}[theorem]{命题}
\newtheorem{definition}[theorem]{定义}
\newtheorem{corollary}[theorem]{推论}
\newtheorem{remark}[theorem]{说明}
\newtheorem{lemma}[theorem]{引理}

% ---- 其他 ----
\usepackage{enumitem}
\usepackage{fancyhdr}
\pagestyle{fancy}
\fancyhf{}
\fancyhead[L]{\small 从Transformer数学结构到Agent工程约束}
\fancyhead[R]{\small arXiv初稿~~\today}
\fancyfoot[C]{\thepage}
\renewcommand{\headrulewidth}{0.4pt}

\title{\textbf{从Transformer数学结构到Agent工程约束\\[10pt]——一个收敛理论框架}}
\date{\today}
\author{}

\begin{document}
\maketitle

\begin{center}
\rule{0.8\textwidth}{0.5pt}
\end{center}

% ============================================================
\begin{abstract}
\noindent
本文论证：大语言模型Agent的一系列已知限制——任务分解的结构性脆弱、记忆缺乏时间动态厚度、全自动闭环的不可达——不是工程上待解决的暂时问题，而是Transformer架构的数学结构所决定的必然结果。

论证分四层展开。第一层（数学基础）：将注意力机制严格表述为FinStoch范畴中的Markov核——单层softmax注意力产出随机矩阵在数学上无争议；多层堆叠中残差连接的加法结构与态射的可乘复合之间存在结构不匹配，本文诚实标注这是整个框架的最薄弱环节。第二层（收敛理论）：提出训练收敛与测试收敛的双模分类，给出收敛的操作化定义（损失稳定、参数更新趋零、行为可预测），并严格分离"收敛保证稳定性"与"经验上稳定解呈现近似函子性质"两个命题。第三层（结构约束）：推导记忆的结构性二分（空间结构vs顺序结构）、上下文驱动的分解决策不可靠的结构必然性（三段论论证）、以及全自动闭环在目标设定层的原则不可闭合性。第四层（工程含义）：导出五项根本设计原则，划定当前架构下的可行域边界。

本文不声称该框架已完备。引言中标注了各节的论证强度分层，相关工作部分逐一标注了与已有文献的差异化路径，讨论部分列出了三条可检验推论及其证伪条件，局限部分诚实标注了当前论证的薄弱环节——包括截断投影命题在有限范畴设定下需要重构、交叉熵不蕴含函子性、以及残差加法与态射复合的结构不匹配。本文的目标是提出一个足以组织现有经验观察的统一理论结构，供社区检验、修正或反驳。
\end{abstract}

% ============================================================
\section{引言}

\subsection{问题}

大语言模型（LLM）在作为Agent运行时，暴露出一系列系统性的、可复现的失败模式：

\begin{itemize}[leftmargin=*]
    \item \textbf{任务分解不可靠。} 模型在多步规划中系统性偏离目标。早期全自动Agent（AutoGPT等）在开放域长程任务中几乎全部失败\cite{AutoGPT2023}；即便是配备了成熟脚手架的生产级Agent（Cursor、Claude Code），仍需在关键决策点设置人工确认。
    \item \textbf{记忆缺乏时间深度。} 模型在长对话中混淆时序、遗忘早期决策、无法区分"三小时前说的"和"刚才说的"\cite{Liu2024LostMiddle}。这些失败无法仅靠扩大上下文窗口消除。
    \item \textbf{全自动闭环失败。} 从目标设定到任务完成的全自动闭环，在开放域中从未被可靠实现；每次"近乎成功"的运行都依赖人类在过程中隐性注入意图澄清和需求仲裁\cite{Anthropic2024Agents}。
\end{itemize}

现有文献将这些失败归类为工程问题——提出更好的分解策略\cite{LEAD2025}、更长的上下文窗口、更优的记忆模块\cite{HAMR2025,CognitiveMemory2025}、更复杂的多Agent协作框架\cite{MACI2025,ChainOfAgents2024}。这些工作在各自的设定下取得了经验改进。

\textbf{本文的立场与此不同。} 我们认为，上述现象不是待解决的工程问题——它们是架构数学结构决定的边界。一个系统在既有架构内的最优工程方案，无法越过其数学结构划定的边界。

\subsection{论证路径}

本文沿以下路线展开：

\begin{enumerate}[leftmargin=*,label=§\arabic*]
    \item \textbf{数学基础}——注意力机制作为FinStoch范畴中的Markov核；多层堆叠中残差加法与态射复合的结构关系；语义范畴的二层构造（有限维底层 + 无限维解释空间）。
    \item \textbf{收敛理论}——收敛的操作化定义；训练收敛与测试收敛的严格区分；交叉熵损失与函子性的分离论证。
    \item \textbf{记忆的结构约束}——空间结构与顺序结构的异质性；无中间过渡的结构论证；时间性的缺失。
    \item \textbf{测试收敛理论}——定义、性质、有效性条件；与训练收敛的系统对比；与Test-time Scaling文献的关系。
    \item \textbf{上下文边界与任务分解}——信息边界的定义；分解的结构不可靠性三段论论证；错误传播机制。
    \item \textbf{层级收敛与不可闭合性}——多级需求收敛结构的形式推导；根需求的外在性；全自动闭环的三层区分。
    \item \textbf{工程含义}——五项根本原则；可行域边界。
    \item \textbf{相关工作、讨论与局限}——与六个子领域文献的系统对话；可检验推论与证伪条件；开放问题与诚实标注。
\end{enumerate}

\subsection{论证强度分层说明}

\textbf{本文各节处于不同的论证层次。} 建议读者根据不同层次的论证强度，选择性地评估和采纳本文的各个部分。

\begin{itemize}[leftmargin=*]
    \item \textbf{严格层次（数学事实）。} §2.1--§2.2中注意力softmax产生Markov核、FinStoch作为工作范畴——这些是数学上无争议的构造，不依赖本文的任何其他假设。
    \item \textbf{形式论证层次。} §3--§9的核心论证链——收敛的操作化定义、记忆的结构性二分、上下文驱动的分解决策不可靠（三段论）、全自动闭环的不可闭合——依赖两个有独立经验证据的前提：(a) 参数空间有效维度远小于名义维度\cite{Li2018IntrinsicDim,Aghajanyan2020,Hu2021LoRA}；(b) 上下文窗口中的信息不经训练收敛筛选。两个前提均有独立于本文框架的经验支持。
    \item \textbf{类比层次的构造。} §2.3中从有限维FinStoch到无限维"语义空间"的扩展，以及米田嵌入在此扩展中的角色——目前是为深层统一性服务的构造，其严格的数学地位有待进一步形式化（§12.3）。
    \item \textbf{经验推论层次。} §10的工程原则和§12.1的可检验推论是从理论框架导出的经验命题，其有效性可独立于理论框架进行检验。如果推论被实验证伪，框架的部分或全部需要修正。
\end{itemize}

\subsection{贡献}

\begin{enumerate}[leftmargin=*,label=\textbf{C\arabic*}.]
    \item 将注意力机制严格表述为FinStoch范畴中的Markov核，在此基础上分析多层堆叠中残差加法与态射复合的结构关系——引入Enriched范畴和Effectus理论作为可能的强闭合路径。
    \item 提出训练收敛与测试收敛的双模分类，严格分离"收敛保证稳定性"与"经验上稳定解呈现近似函子性质"两个命题。
    \item 从空间结构与顺序结构的数学异质性出发，论证LLM记忆中不存在长短期之间的连续过渡。
    \item 给出上下文驱动的任务分解不可靠的结构论证（三段论形式）。
    \item 论证全自动Agent闭环在目标设定层的原则不可闭合性，并给出三层区分。
\end{enumerate}


% ============================================================
\section{理论基础：注意力作为FinStoch中的Markov核}

\subsection{工作范畴：FinStoch}

本文选择的工作范畴是\textbf{FinStoch}——有限集上的随机映射范畴。选择该范畴的理由是：(a) 其数学结构已被充分刻画\cite{Fritz2020Markov}；(b) Transformer注意力的核心计算精确地产生该范畴中的合法态射。

\begin{definition}[FinStoch范畴]
    FinStoch的构造如下：
    \begin{itemize}[leftmargin=*]
        \item \textbf{对象：} 有限集。在语言模型中，退化为词汇表$V$上的概率单纯形$\Delta^{|V|-1}$。
        \item \textbf{态射：} 随机矩阵（Markov核）。$f: A \to B$是一个$|A| \times |B|$矩阵，每行和为1，即条件概率分布$P(B \mid A)$。
        \item \textbf{复合：} 矩阵乘法——对应Chapman-Kolmogorov积分$\int P(C \mid B)P(B \mid A)\,dB$。结合性由矩阵乘法结合律保证。
        \item \textbf{恒等态射：} 单位矩阵$I$（对应Dirac delta $\delta(A' \mid A)$的离散化）。
    \end{itemize}
\end{definition}

在此范畴中，说"语义间贝叶斯关系构成态射"不是隐喻——它是数学事实。FinStoch的态射精确地编码了条件概率结构，而条件概率结构正是分布语义学的数学模型。

\begin{remark}[与米田引理的关系]
    FinStoch中一个对象$X$的"语义身份"由所有其他对象到$X$的条件概率（即所有以$X$为目标的态射）所唯一确定。这恰好是分布语义学（Firth 1957："you shall know a word by the company it keeps"）的范畴论精确表述。Bradley, Terilla \& Vlassopoulos (2022)\cite{Bradley2022Enriched}在$[0,1]$-enriched范畴中显式地使用米田嵌入定义了语义——"文本的意义是其representable copresheaf"——将分布语义假设提升为范畴论定理。米田嵌入$Y: \mathcal{C} \to [\mathcal{C}^{\mathrm{op}}, \mathbf{Set}]$提供了该直觉的数学定理：对象由指向它的所有态射完全确定。在FinStoch中，我们以随机矩阵而非预层为态射——该选择使注意力机制的直接数学对应成为可能，同时保留了分布语义学的核心直觉。
\end{remark}

\subsection{单层注意力严格是Markov核}

单头自注意力的核心计算为：

\begin{equation}
    \mathrm{Attention}(Q,K,V) = \mathrm{softmax}\!\left(\frac{QK^{\mathsf{T}}}{\sqrt{d_k}}\right)V
    \label{eq:attention}
\end{equation}

令$A = \mathrm{softmax}(QK^{\mathsf{T}}/\sqrt{d_k})$。$A$是一个$n \times n$矩阵（$n$为序列长度），其中$A_{ij} \in [0,1]$且$\sum_{j=1}^{n} A_{ij} = 1$（每行经softmax归一化）。

\begin{proposition}[注意力矩阵是Markov核]
    $A$在数学上\textbf{严格是一个随机矩阵}，即FinStoch中的一个合法态射。$A_{ij}$编码了条件概率$P(\text{位置}j\text{被关注} \mid \text{位置}i\text{发出query})$。

\textbf{多头注意力的结构。} 每个注意力头$h$独立计算其注意力矩阵$A^h = \mathrm{softmax}(Q_h K_h^{\mathsf{T}}/\sqrt{d_h})$，每个$A^h$均为合法Markov核。各头用$A^h$对值向量的不同子空间加权（$\mathrm{head}_h = A^h \cdot V W_h^V$），所有头输出被拼接后经$W^O$做线性组合。关键点：$W^O$组合的是各头在\textbf{值空间中的输出向量}，而非注意力矩阵本身。$A^1, A^2, \dots, A^H$从未被合并为单一随机矩阵——多头结构是$H$个\textbf{并行独立}的Markov核的集合，而非单一"凸组合"核。

\textbf{这不影响后续论证。} 后续全部推导（记忆二分、测试收敛、分解决策不可靠、不可闭合）依赖的前提仅为"单层注意力产生合法的Markov核"——该前提对每个头独立成立。多头结构不改变注意力作为FinStoch态射的数学地位。
\end{proposition}

\textbf{这一步没有近似，没有比喻。} 这是整个框架中论证强度最高的环节。Santos et al. (2026)\cite{Santos2026Sparse}提供了严格的形式证明：softmax注意力与高斯核Nadaraya-Watson估计器之间的精确函数等价。Chen \& Li (2023)\cite{Chen2023SGPA}进一步证明了当点积被替换为对称核时，注意力输出精确等于稀疏变分高斯过程（SVGP）后验均值。Kang, Lee \& Cheng (2026)\cite{Kang2026PPD}从构造上证明transformer可通过Richardson迭代在其前向传播中精确计算GP后验均值和方差。

注意力作为Markov核的数学地位也直接连接了它与Nadaraya-Watson非参数核估计。当核函数选为$\exp(QK^{\mathsf{T}}/\sqrt{d_k})$时，Nadaraya-Watson估计器$\hat{f}(q) = \sum_i K(q, k_i)y_i / \sum_j K(q, k_j)$的输出精确等于softmax注意力的输出\cite{Santos2026Sparse}。在高斯过程先验下，该估计器产生后验均值——但此处的精确关系是：Nadaraya-Watson估计器等价于\textbf{忽略训练点间协方差后的简化GP后验均值}。它不是完整的GP推断——该区分对于避免过度声称很重要。

\subsection{多层堆叠与态射复合的结构关系}

这是全论证中\textbf{最关键的环节}，也是本文需要诚实标注的薄弱点。

\textbf{理论要求（态射复合）：} 如果第1层注意力核为$A_1$，第2层为$A_2$，则两层的"复合态射"应为$A_{\mathrm{eff}} = A_2 \cdot A_1$（矩阵乘法 = Chapman-Kolmogorov积分）。

\textbf{实际计算（Transformer层）：}
\begin{align}
    h_l' &= x_l + \mathrm{Attn}_l(\mathrm{LN}(x_l)) \label{eq:transformer1}\\
    h_l  &= h_l' + \mathrm{MLP}_l(\mathrm{LN}(h_l')) \label{eq:transformer2}\\
    x_{l+1} &= h_l \label{eq:transformer3}
\end{align}

其中$\mathrm{Attn}_l(\mathrm{LN}(x_l)) = A_l \cdot V_l \cdot W_l^O$。

展开$L$层后的残差流：
\begin{equation}
    x_L = x_0 + \sum_{l=0}^{L-1} \bigl[A_l \cdot V_l \cdot W_l^O + \mathrm{MLP}_l(\cdots)\bigr]
    \label{eq:residual}
\end{equation}

\textbf{关键观察：} 残差流中的信息是各层输出的\textbf{和}，而非各层注意力核的\textbf{积}。$A_l$和$A_{l-1}$之间不存在直接的矩阵乘法关系。

\begin{table}[htbp]
\centering
\caption{态射复合与Transformer实际行为的结构对比}
\label{tab:composition_gap}
\begin{tabular}{@{}p{2.5cm}ll@{}}
\toprule
\textbf{维度} & \textbf{态射复合（理论）} & \textbf{Transformer（实际）} \\
\midrule
组合方式  & 乘法：$A_2 \cdot A_1$ & 加法累积：$x_0 + \sum F_l(x_l)$ \\
核间关系 & 直接复合，C-K积分 & 通过残差流中间接关联 \\
结构保持 & 严格保持 & 功能逼近——深层表示逐渐编码多步依赖 \\
\bottomrule
\end{tabular}
\end{table}

\begin{remark}[为什么功能上仍然可能工作]
    虽然$A_l$不直接相乘，但残差流提供了另一种累积信息的方式：(1) 短期依赖由$A_l$直接捕获；(2) 中期依赖通过残差累积建立——第$l+1$层处理的$x_l$已包含第$l$层的信息，$A_{l+1}$可据此分配注意力；(3) 长期依赖由多个残差块叠加产生。Transformer不是在用矩阵乘法复合Markov核，而是在用残差加法迭代地修正一个逐渐包含更多上下文信息的表示。训练中的梯度引导模型在残差流中安排信息，使$A_{l+1}$作用在包含$A_l$效果的表示上的行为，在功能上逼近$A_{l+1} \cdot A_l$。
\end{remark}

\textbf{闭合状态：} 单层注意力$=$Markov核——严格闭合。多层堆叠$=$态射复合——仅可弱闭合（功能逼近），因为残差加法与可乘复合之间存在结构不匹配。强闭合需要引入更复杂的范畴论工具（如Effectus理论\cite{Cho2015Effectus}中的加性操作，或Para构造\cite{Fong2019Para}中的参数化态射双层结构），当前框架仅标注这些路径而未完成其形式化。这是本文理论构造中最薄弱的环节。

\subsection{语义范畴的二层构造}

本文的框架需要在两种"语义空间"之间建立区分，以解决§12.3中讨论的截断投影命题的设定问题。如果工作范畴仅为FinStoch（有限对象集），则语义空间本身是有限维的，"截断"作为原则性限制不成立。

\textbf{二层构造：}
\begin{itemize}[leftmargin=*]
    \item \textbf{底层（实现层）：} FinStoch中的token分布间Markov核——$|V|$维，有限。这是模型直接工作于其上的范畴。不需要截断。
    \item \textbf{上层（解释层）：} 连续语义空间——所有可能的概念及其语义关系的完备空间。该空间的"完整米田嵌入"（如果形式化地存在）需要无限维表征，因为概念的数量和关系的复杂性在原则上无界。模型学到的FinStoch中的态射是该完整空间在有限参数流形上的\textbf{投影}。
\end{itemize}

在此二层构造下，"截断投影"的精确含义是：从无限维的解释语义空间到有限维模型表征的映射过程中必然存在信息压缩。该构造提供了深层统一的直觉框架，但其严格的数学定义——特别是上层空间的范畴论构造及其与底层FinStoch之间的映射——是开放问题（§12.3）。


% ============================================================
\section{收敛理论}

\subsection{收敛的操作化定义}

"收敛"在优化理论中有狭窄的技术含义（驻点性、全局最优性条件等）。本文在更基础的层面使用该词——指系统经过迭代优化达到稳定状态。为避免歧义，本节给出操作化定义，后续全篇一致使用。

\begin{definition}[收敛（本文用法）]
    模型参数在梯度下降的迭代驱动下，以训练数据上的损失函数为引导，逐步进入稳定配置的过程。完成的标志是三项可观测条件：
    \begin{enumerate}[leftmargin=*]
        \item 损失曲线停止显著下降（训练loss进入平台区）；
        \item 参数更新幅度趋近于零（梯度范数/权重更新量低于阈值）；
        \item 模型行为在训练分布上转入可预测的稳定模式（同分布样本下的输出分布不再显著漂移）。
    \end{enumerate}
\end{definition}

本文说一组映射"受收敛约束"，仅指它经历了上述迭代稳定化过程。不附加全局最优性、唯一性等更强的数学承诺。

这个界定的核心意义在于：\textbf{不是所有进入模型的信息都经过了收敛。} 有些信息（参数）进入了收敛的熔炉，在损失函数的持续压力下被反复锻造；有些信息（上下文）直接注入模型的前向计算，从未经历过一步梯度反馈。这一区分是本文记忆理论（§4）和测试收敛理论（§5）的基石。

\subsection{训练收敛}

\begin{definition}[训练收敛]
    以训练数据上的损失函数（通常为交叉熵$L = -\sum \log P(\mathrm{next\_token} \mid \mathrm{context})$）为驱动信号，通过梯度下降在参数空间中完成的收敛过程。收敛完成后参数静态冻结——模型转入只读模式。
\end{definition}

训练收敛的\textbf{范围}由两重约束共同决定：
\begin{enumerate}[leftmargin=*]
    \item \textbf{分布约束：} 仅在训练数据中实际出现的语义关系的分布范围内提供稳定性保证。训练数据中没有标注的语义关系——无论它们在客观语义空间中是否成立——不在收敛覆盖范围内。
    \item \textbf{维度约束：} 参数空间的有效维度（本征维度）远小于名义维度。Li et al. (2018)\cite{Li2018IntrinsicDim}发现MNIST FC网络仅需$\sim$750维（总参数的0.4\%）即可达到满性能。Aghajanyan et al. (2020)\cite{Aghajanyan2020}在语言模型\textbf{微调}场景中发现：RoBERTa-Large（3.55亿参数）在MRPC任务上仅需$\sim$200个可训练参数即可达到满微调性能的90\%——压缩比约170万倍。LoRA\cite{Hu2021LoRA}以$r=1$--$16$即可有效适配大模型（同样在微调场景）。Mabrok (2026)\cite{Mabrok2026Latent}发现在六种Transformer架构的\textbf{预训练表示}中，隐藏层本征维度仅占环境空间的1--3\%。Xu (2026)\cite{Xu2026Holographic}发现\textbf{预训练grokking轨迹}被限制在2--6维全局子空间。Kalyoncuoglu \& Miftachov (2026)发现表示可用随机投影压缩16倍而准确率仅损失$\sim$1\%。

\textbf{外推警示：} 从微调场景（Aghajanyan 2020, LoRA）到预训练表征的"语义空间本身低维"之间存在非平凡外推。微调的低秩有效性证明的是\textbf{任务适配}只需低维子空间，不等价于预训练表征的全部语义内容均可压缩至低维。Mabrok (2026)和Xu (2026)为预训练阶段的低维性提供了更直接的证据，但这些研究较新，尚需更多独立验证。本文的核心论证依赖的是"训练收敛的\textbf{有效范围}有限"——即使参数空间的本征维度随模型规模有某种程度的增长，收敛约束的边界仍然存在（该边界由训练数据分布的有限性决定，而非仅由参数空间维度决定）。
\end{enumerate}

\subsection{收敛与函子性的分离}

\textbf{这是本文框架的一个关键分离。}

监督训练的损失函数是交叉熵：
\begin{equation}
    L = -\sum \log P(\mathrm{next\_token} \mid \mathrm{context})
\end{equation}
该损失促使模型匹配训练数据中的条件分布。它\textbf{不}促使模型满足函子性条件$\mathcal{F}(g \circ f) = \mathcal{F}(g) \circ \mathcal{F}(f)$。交叉熵和函子性条件在数学上没有包含关系——前者不以任何形式蕴含后者。

\begin{proposition}[收敛与函子性的分离]
    收敛保证的是\textbf{稳定性}（参数配置在训练分布上不再漂移）。经验上，在自然语言数据上收敛到的稳定配置呈现出近似的函子性质——模型在训练分布内保持了语义关系的复合结构。但这不是收敛定义的逻辑蕴含，而是数据本身的统计结构（语义关系在自然语言中确实具有近似函子的复合性）和模型架构的归纳偏置共同导致的\textbf{涌现效果}。
\end{proposition}

\textbf{这个分离影响了全篇的论证承载：} 当我们说模型在截断投影范围内具有"结构保持的映射"时，"结构保持"是一种经验发现而非逻辑后承。它的可靠性在分布内足够高以至于工程上可用，但在分布边界附近缺乏理论保证。这恰好是§5--§7中论证"上下文驱动的分解决策不可靠"的数学根源——分解决策位于或超出分布边界时，经验上涌现的函子性不再适用。

\subsection{思维链RL后训练的作用}

CoT RL后训练的有效性已被多项工作证实\cite{OpenAI2024o1,Kaiser2025FirstPrinciples}。本文的分析如下：

\begin{proposition}[CoT RL不扩展收敛的语义边界]
    CoT RL后训练改变的是优化目标（从逐token匹配到多步推理一致性），但参数仍然在同一个本征维度受限的参数流形上更新。RL的作用是在截断投影\textbf{内部}重新组织映射结构——"拉直"被截断投影压缩后不一致的态射复合路径——使得原先隐式可推导但不稳定的中间推理步骤变得显式化且可复现。
\end{proposition}

\textbf{这解决了原框架中的一个张力：} 如果低秩空间装不下逻辑结构（§2.4的维度约束），CoT RL如何能"补全"逻辑能力？答案是：低秩空间\textbf{装得下}逻辑结构——预训练之所以没有良好地表征它，不是因为容量不够，而是因为预训练目标（逐token条件概率匹配）没有引导参数去组织逻辑结构的表征。CoT RL改变了优化目标，使参数在已有容量内重新组织了流形。因此"维度有限"和"CoT RL有效"可以同时成立——前者不构成后者的逻辑障碍，后者不否定前者作为分布边界的存在。


% ============================================================
\section{记忆的结构性二分}

\subsection{两种异质存储}

模型运行时可访问的全部信息来自两个来源。这两个来源在数学结构上是\textbf{异质的}——不仅是容量或持久性上的差异。

\begin{table}[htbp]
\centering
\caption{LLM中两种存储机制的结构对比}
\label{tab:memory}
\begin{tabular}{@{}p{2.8cm}ll@{}}
\toprule
& \textbf{长期（参数）记忆} & \textbf{短期（上下文）记忆} \\
\midrule
载体      & 模型权重矩阵 & 上下文窗口中的token序列 \\
形成方式  & 训练收敛——在损失压力下反复迭代 & 每次推理时顺序注入 \\
数学结构  & 参数空间中的点（\textbf{空间结构}） & 序列中的线性展开（\textbf{顺序结构}） \\
收敛约束  & 经受 & 未经受 \\
动态性    & 无（训练后冻结） & 有顺序性，无时间性 \\
持久性    & 跨会话持久 & 会话内挥发性 \\
\bottomrule
\end{tabular}
\end{table}

关键差异在\textbf{数学结构}层面：参数记忆是空间结构——信息以分布式表征的形式被压缩在整个权重矩阵中，各个事实不是独立条目，而是权重空间中的一个点。上下文记忆是顺序结构——信息以线性序列逐token注入，位置由位置编码标记。

\subsection{无中间过渡的结构论证}

2025年的多项工作提议在长期参数记忆和短期上下文记忆之间增设"中期记忆层"\cite{HAMR2025,MemGPT2024,MemoryBear2025}。我们的论证是：这些方案并未创建新的记忆\textbf{结构}，而是在管理短期记忆的\textbf{容量}。

\begin{proposition}[无中间过渡]
    空间结构（参数空间中的分布表征点）与顺序结构（序列中的线性排列）属于不同的数学范畴。从顺序到空间不存在连续的、双向可逆的映射。具体地：
    \begin{enumerate}[leftmargin=*]
        \item 一个序列可被压缩为固定长度向量（摘要嵌入）→ 该向量进入了空间结构。它不再具有原序列的顺序可遍历性。
        \item 空间中的一个点可被"展开"为文本 → 该文本重新进入了顺序结构。
        \item 不存在一个\textit{既具有空间结构的持久性又具有顺序结构的动态可遍历性}的独立中间层。
    \end{enumerate}
\end{proposition}

所谓的"中期记忆"方案本质上是两种操作的交替：(a) 将顺序信息压缩进空间结构（摘要→嵌入→存储），此时它成为静态的长期记忆；(b) 从存储中检索片段重新注入上下文窗口（检索→展开→注入），此时它被拉回短期记忆。没有独立的"中间"数学结构。这些方案\textbf{有效}——摘要压缩帮助容纳更多历史、向量检索帮助找到相关信息——但它们有效的方式是管理短期记忆的容量（控制"什么进入有限窗口"），而非创建新的记忆结构。

\subsection{时间性的缺失}

\textbf{声称的精确范围。} 本节论证Transformer语言模型中缺乏\textbf{跨会话的时间性动态}——即记忆强度随真实时间流逝而衰减、在离线期间被自发巩固和重组的能力。该声称的否定范围不包含上下文窗口\textbf{内部}的顺序效应（近因效应、首因效应等），后者在Transformer中已被实验证实存在\cite{Mistry2025Episodic,Park2025Temporal}，本文将其归类为"窗口内的顺序性"而非"跨会话的时间性"。窗口内的顺序效应和跨会话的时间性动态之间的区分在§12.3中进一步讨论。

\begin{definition}[顺序 vs 时间性]
    位置编码提供的是\textbf{顺序}（$A$在$B$之前），不是\textbf{时间性}（$A$是三小时前说的，$B$是刚才说的）。顺序是对称的——$A$在$B$前等价于$B$在$A$后。时间性是不对称的——"三小时前"和"一分钟前"在提取难易度、记忆保真度、对当前行为的影响权重上不同。
\end{definition}

人类情景记忆的时间性维度——Ebbinghaus遗忘曲线（随时间衰减）、睡眠巩固（离线重组）、情感调制、干扰效应、情境依赖提取——在认知科学中通常被视为情景记忆的核心特征。这些特征在Transformer中缺失：所有位置的token在注意力机制下平等可及——无衰减、无巩固、无自发重组。这不是窗口大小的问题——即使把窗口扩到1亿token，信息仍按顺序排列，无时间厚度。

\begin{corollary}[工程后果]
    所有在现有Transformer架构内试图实现"渐进式持久记忆"的方案，面临一个共同的结构限制：它们可以更高效地管理上下文窗口的\textbf{容量}，但无法改变窗口内信息的\textbf{性质}——它仍然是顺序性的，没有时间厚度。Hu et al. (2025)\cite{Hu2025AgentMemory}在一项覆盖200+论文的全面综述中独立确认了这一区分：\"Agent记忆\"需要跨交互的持久、演化、自我改进的知识存储，而现有方案主要提供静态检索——缺乏形成$\to$演化$\to$检索的完整生命周期。Shan et al. (2025)\cite{Shan2025Memo}从互补学习系统（CLS）理论出发，独立论证了当前Agent记忆是\"备忘录而非真正记忆\"——海马体（快速情景存储）已被实现，但新皮层（缓慢的权重巩固/抽象）缺失。
\end{corollary}

\textbf{反证与讨论。} 我们诚实地注意到以下发现对\"时间性缺失\"论断构成挑战：Park et al. (2025)在Transformer中发现了特定的\"时间头\"（Temporal Heads）——专门处理时间知识的注意力头，禁用它们会选择性损害时间相关回忆；Mistry et al. (2025)发现GPT-2注意力头展示了\textbf{时间邻近性、首因效应和近因效应}——人类情景记忆的核心成分；Wang et al. (2025)的Token-Time假说实验表明LLM具有可测量的\"时间共情\"——将离散token数映射到连续墙钟时间。这些发现表明\textbf{有限的时间感知}在Transformer表示中确实自发涌现了。但关键区分是：这些效应是在单次推理的上下文窗口内发生的顺序效应，而非跨会话的、随时间自发衰减和巩固的时间性动态。\"时间头\"和\"近因效应\"是\textbf{窗口内的位置敏感性}——$A$在$B$之前且模型倾向于关注$A$更近——而非Ebbinghaus式的、随真实时间流逝而变化的记忆强度曲线。本文的论断\"无时间性\"的精确范围是\textbf{跨会话的时间性动态}，而非窗口内的顺序效应。两者的区分在§12.3中进一步讨论。


% ============================================================
\section{测试收敛：Agent行为的支配原理}

\subsection{为什么训练收敛不能解释Agent行为}

训练收敛（§3.2）描述模型\textbf{能力形成}的过程。Agent行为涉及的是一个不同的问题：在参数冻结的条件下，如何在上下文窗口中通过多轮交互逐步逼近目标。两者的认知性质不同：

\begin{itemize}[leftmargin=*]
    \item 训练收敛解释"模型\textbf{有能力}做什么"——它在训练数据中学会了哪些局部语义态射。
    \item 训练收敛不解释"模型\textbf{在这一次}是否完成了一个需要多步复合态射的任务"——复合后的一致性是不由局部态射保持性自动保证的全局属性。
\end{itemize}

这解释了为什么训练得再好的模型，作为Agent时在多步任务中仍然会偏离——驱动Agent行为的认知过程与驱动模型能力形成的过程具有\textbf{不同的数学保证级别}。

\subsection{测试收敛的定义}

\begin{definition}[测试收敛]
    设Agent在上下文窗口内执行任务。令$\mathcal{T}$为目标状态，$\mathcal{V}$为一组可自动执行的验证操作。Agent通过循环
    \[
    \text{生成} \longrightarrow \text{验证} \longrightarrow \text{分析差距} \longrightarrow \text{修正}
    \]
    逐步逼近$\mathcal{T}$。这一过程称为\textbf{测试收敛}。
\end{definition}

\begin{table}[htbp]
\centering
\caption{训练收敛与测试收敛的系统对比}
\label{tab:two_convergences}
\begin{tabular}{@{}p{3cm}ll@{}}
\toprule
& \textbf{训练收敛} & \textbf{测试收敛} \\
\midrule
驱动信号      & 监督损失（连续梯度向量场） & 可验证目标（离散通过/失败） \\
收敛保证      & 有（梯度下降的数学性质） & 无（依赖验证信号质量与覆盖度） \\
作用范围      & 参数空间（$\Delta W$） & 上下文空间（$\Delta \text{output}$） \\
持久性        & 静态（训练后冻结） & 挥发性（过程结束即丢弃） \\
局部最优逃逸  & 随机梯度噪声+SGD动量 & 无自动逃逸机制 \\
\bottomrule
\end{tabular}
\end{table}

\subsection{测试收敛的有效性条件与失效模式}

测试收敛的有效性完全来自验证信号$\mathcal{V}$的覆盖度和精度。设任务目标$\mathcal{T}$的约束集合为$\mathcal{C} = \{c_1, \dots, c_n\}$，验证信号覆盖子集$\mathcal{C}_{\mathcal{V}} \subseteq \mathcal{C}$。

当$\mathcal{C}_{\mathcal{V}} = \mathcal{C}$时，测试收敛在原理上可趋近$\mathcal{T}$。当$\mathcal{C}_{\mathcal{V}} \subsetneq \mathcal{C}$时，Agent可通过所有显式测试但最终输出不满足实际需求——\textbf{这不是模型能力不足，是验证信号的信息量不足以唯一确定目标状态。}

\textbf{失效模式：} 2025年的独立实证研究系统性地记录了测试收敛的失效模式。Gema et al.\cite{Gema2025Inverse}发现了五种与本文框架预测一致的"逆缩放"模式——模型在推理时思考得更久反而更差（分心、伪相关回归、逻辑过度复杂化、过度思考）。Wang et al.\cite{Wang2025TTSPM}提出了Test-Time Scaling Performance Model（TTSPM）量化了收益递减的精确数学形式：$P(C) \approx P_0 + \alpha \log C$。Zhao et al.\cite{Zhao2025Verification}的实验确认了验证是瓶颈而非生成——模型在孤立状态下难以发现自身错误，但在跨样本比较时可以定位错误位置。

这些发现共同支持了一个核心判断：\textbf{测试收敛缺乏训练收敛的数学保证，这不是工程缺陷，而是两种驱动机制之间的根本差异。}

\subsection{与Test-time Scaling文献的关系}

现有Test-time Scaling文献\cite{SETS2025,IAD2025,ReVeal2025}研究"推理时多算几步如何提升性能"。本文追问的是不同的问题：\textbf{推理时多算几步为什么有时候有效有时候无效？}

本文的回答：有效与否取决于验证信号是否足够密集地覆盖目标结构的约束。数学证明（验证可自动执行）→ 测试收敛工作条件良好 → test-time compute稳定提升。创意写作（"好"无法被自动测试完全捕获）→ 测试收敛缺失 → 增加compute无法替代验证信号的消歧作用。OpenReview 2025的一篇论文给出了与本文方向一致的形式证明：verifier-based方法在test-time预算增长时渐进地支配verifier-free方法，差距约按$\sqrt{H}$扩大——验证信号的存在与否决定了推理时扩展的有效性上界。


% ============================================================
\section{上下文边界与任务分解}

\subsection{信息边界}

对于单次连贯推理，模型可获取的全部信息来自：参数中的长期记忆（静态）+ 上下文中的短期记忆（顺序注入）。上下文窗口的物理长度（ContextCap）不等于\textbf{有效}利用长度（EffCap）——Lost in the Middle效应\cite{Liu2024LostMiddle}已系统验证。

\begin{definition}[信息边界]
    $\mathrm{EffCap}$为模型能从上下文中可靠提取结构保持映射的最大信息量。已知$\mathrm{EffCap} \ll \mathrm{ContextCap}$。超出信息边界的任务——部分约束落在模型有效注意范围之外——无法被单次连贯处理完成。
\end{definition}

\subsection{分解是为了恢复测试收敛的工作条件}

超出信息边界的任务\textbf{必须}分解。但分解的\textbf{目的}需要精确表述：不是为了提高性能，而是为了\textbf{恢复测试收敛的工作条件}——让每个子任务拥有独立的上下文窗口，在该窗口内信息量不超出EffCap，使测试收敛有运作空间。

\subsection{分解行为本身也受边界约束}

分解行为——判断是否分解、如何分解、接口定义——发生在上下文窗口内，依赖上下文中的信息。但上下文中的信息不受收敛约束（§3.2、§4.1）。

\textbf{自指约束：} 如果任务信息量没有超出边界，不需要分解；如果需要分解，分解决策所需的信息量可能已经超出了边界。这是一个信息论上的良性循环约束——不是"有时会失败"的经验巧合，而是从前提推导出的结构结论。


% ============================================================
\section{任务分解的结构不可靠性}

\subsection{错误沿层级向下传导}

设节点$P$分解为子节点$\{S_1, S_2, \dots, S_k\}$。若$P$的指定存在语义偏差（$P$不能唯一收敛到用户真实意图），所有$\{S_i\}$以该偏差为基础展开。$S_i$只持有$P$赋予的局部上下文，不具独立检测和纠正$P$层级错位所需的全域信息。这是通信理论中错误传播的语义版本——但更严重，因为传播的不是比特错误而是语义坐标的位移。

\subsection{三段论：上下文驱动分解不可靠的结构必然性}

\begin{enumerate}[leftmargin=*]
    \item \textbf{大前提：} 上下文环境中的信息不经收敛约束（§3.2、§4.1：上下文绕过训练收敛筛选，其中的信息没有经过参数空间的结构压缩和损失函数的反馈锻造）。
    \item \textbf{小前提：} 任务分解需要模型在上下文窗口内做出需要可靠复合态射保持的判断——识别全局依赖、判断子任务耦合度、定义接口边界。
    \item \textbf{结论：} 用不经收敛约束的信息来做出需要收敛级可靠性的决策——不可靠是\textbf{逻辑上的结构必然}，而非经验上的偶然。
\end{enumerate}

\begin{corollary}
    将任务分解的结构决策（架构选择、接口定义、语义分歧消解）交由模型在上下文中自主完成——这种设计的失败率有理论下界。该下界不由prompt质量决定，而由"上下文中的信息不经收敛约束"这一前提决定。
\end{corollary}

\subsection{与经验文献的归因层差异}

LEAD\cite{LEAD2025}、Divide-and-Conquer噪声模型、CoDA等2025年文献从经验层面确认了分解脆弱性，归因为噪声积累、规划-执行脱节、上下文过载。本文的归因更深一层：\textbf{上下文本身不受收敛约束}。噪声归因暗示"去噪即可接近可靠"；结构归因指出：即使零噪声，用无约束的信息做需约束的决策，逻辑上就不保证可靠性——噪声降低可减小失败的方差，不能消除失败的结构根源。

\subsection{关键分叉点须外部介入}

不可逆的高影响决策节点（需求消歧、顶层接口定义、依赖关系声明）必须由\textbf{系统外部}——人工操作者，或独立于模型上下文窗口的硬约束程序（编译器、类型检查器、自动化测试套件）——进行仲裁。这对应了当前最成功的Agent系统（Cursor、Claude Code）的设计模式，也解释了为什么全自动早期Agent（AutoGPT）在开放域中几乎全部失败。


% ============================================================
\section{层级收敛与不可闭合性}

\subsection{多级需求收敛结构的形式推导}

本节给出§5--§7结论的形式化推导。前提和桥接命题如下：

\begin{table}[htbp]
\centering
\caption{推导前提}
\label{tab:premises}
\begin{tabular}{@{}ccl@{}}
\toprule
\textbf{编号} & \textbf{类型} & \textbf{陈述} \\
\midrule
P1 & 定义 & 信息边界 = 模型对单一问题的可靠处理能力边界 \\
P2 & 推论 & 超出边界的问题不受收敛覆盖（§3.2） \\
P3 & 定义 & Agent通过测试收敛逼近目标——可验证目标+反复测试（§5.2） \\
P4 & 分析 & 测试收敛依赖可验证目标的存在（§5.3） \\
P5 & 推论 & 上下文中的信息不经收敛约束（§3.2, §4.1） \\
P6 & 定义 & Agent = Model + Harness（§9.3） \\
\midrule
B1 & 经验 & 多数有意义的任务超出信息边界，分解是必然的 \\
B2 & 结构 & 流程节点是一次性向下传导——节点错误乘性放大 \\
B3 & 概念 & 验证权威递归上溯最终不可自举——终止于需求提出者 \\
\bottomrule
\end{tabular}
\end{table}

\begin{proposition}[推导一：任务必须分解]
    由B1（多数任务超出边界）和P1+P2（超边界不受收敛覆盖），得：任务必须分解为落入边界内的子任务。
\end{proposition}

\begin{proposition}[推导二：子任务须可验证，分解形成层级]
    子任务需可验证目标以触发测试收敛（P3+P4）。分解行为本身在上下文窗口内完成（P1），需测试收敛（P3）。上层收敛产出下层目标——层级自然形成。
\end{proposition}

\begin{proposition}[推导三：硬节点约束]
    上下文映射不经收敛（P5）。流程节点错误代价乘性放大（B2）。因此纯上下文驱动的流程节点管理不可靠——关键分叉点须由经独立验证的逻辑固定（人工或程序硬约束，P6）。
\end{proposition}

\begin{proposition}[推导四：根需求人工锚定]
    测试收敛依赖可验证目标（P4）。验证标准递归上溯——每层目标需被上层验证。最上层验证标准的判断权威不可由系统替代（B3）。递归终止于人工给定的最高层需求。
\end{proposition}

\subsection{全自动闭环的三层区分}

\begin{table}[htbp]
\centering
\caption{三层不可闭合性}
\label{tab:three_layers}
\begin{tabular}{@{}p{2.5cm}p{4cm}p{3cm}p{2cm}@{}}
\toprule
\textbf{层面} & \textbf{定义} & \textbf{不可闭合的性质} & \textbf{论证强度} \\
\midrule
目标设定层 & 系统能否自主生成最高目标 & 逻辑必然不可自举 & \textbf{强} \\
\multicolumn{4}{@{}l}{\small 目标语义不在训练分布内，无从收敛（§3.2, §8.4）} \\
\midrule
任务分解层 & 给定目标下自主规划子任务 & 结构不可靠 & \textbf{中} \\
\multicolumn{4}{@{}l}{\small 上下文驱动分解不受收敛约束（§7.2三段论）} \\
\midrule
执行闭环层 & 给定验证标准下自主完成任务 & 条件可行 & \textbf{弱} \\
\multicolumn{4}{@{}l}{\small 若验证信号覆盖充分，执行可在无人工介入下完成（§5.3）} \\
\bottomrule
\end{tabular}
\end{table}

"\textbf{全自动闭环在原则上不可闭合}"的精确含义是：从目标设定到目标达成的完整链路中，\textbf{最高目标不可由系统自举}——根需求的语义落在收敛空间的结构边界之外。这与"执行过程不能自动化"是不同的主张——在验证信号覆盖充分的条件下，执行闭环层是可以闭合的。

\subsection{与现有AGI不可能论证的路径差异}

已有文献从多条路径论证了AI的原则限制\cite{Schlereth2025ICB,Bui2025Foundations}：哥德尔不完备、图灵停机、中文屋、热力学。本文的路径不同：(1) 起点是可测量经验量（参数空间有效维度有限），而非哲学假设；(2) 结论是可证伪的（如果实验发现本征维度可无限扩展，部分结论需修正）；(3) 论证链仅依赖架构数学结构，不依赖"智能"或"理解"的哲学定义。


% ============================================================
\section{工程含义}

\subsection{Agent = Model + Harness}

\begin{definition}
    $\mathrm{Agent} = \mathrm{Model} + \mathrm{Harness}$，其中Harness = 程序框架 + 上下文环境。模型提供截断的语义映射能力（训练收敛的产物）；Harness提供外部验证、硬约束和状态管理（测试收敛的实现框架）。
\end{definition}

Harness不是"模型能力不够时的权宜之计"。在模型的结构边界不会因能力提升而消失的前提下（边界来自架构数学结构，而非训练不充分），\textbf{Harness的设计是第一性的，不是过渡性的。}

\subsection{五项根本工程原则}

\begin{enumerate}[leftmargin=*,label=\textbf{P\arabic*}.]
    \item \textbf{精确性优先：} 提示词的明确指示——最大限度消除语义歧义维度。精确不等同于短——CoT增加token数但通过消歧降低信息结构复杂度。
    \item \textbf{结构简化：} 最小化输入的歧义维度数量，而非token数。每增加一个可被多重解释的信息维度，输入在语义空间中的可能落点指数增长。
    \item \textbf{语义氛围：} 将输入置于与目标态射最近的训练分布区域——利用训练收敛在分布内的高密度覆盖。
    \item \textbf{分叉控制：} 关键决策节点（任务划分、接口定义、需求消歧）必须由外部仲裁介入——不信任上下文驱动的结构决策（§7.2三段论）。
    \item \textbf{可验证性：} 分解后的每个任务单元必须配备可自动执行的验证标准——Agent行为由测试收敛驱动，测试收敛依赖验证信号（§5.3）。
\end{enumerate}

\subsection{可行域边界}

\textbf{在当前数学结构下可行：} 任务切至EffCap内使测试收敛有条件运作；每单元配备可自动验证标准；结构决策由外部仲裁介入；接受根需求的外在性——人机协作不是过渡方案而是结构必然。

\textbf{在当前数学结构下不可行：} 模型自主规划长程无仲裁任务；模糊需求驱动多步自主决策而无自动验证；在现有Transformer架构内实现完整闭环自举。这不是"技术上困难"——是给定架构数学结构下收敛条件在原则上无法满足。


% ============================================================
\section{相关工作}

本节逐一标注本文与已有文献的差异化路径。差异化不等于否定——已有文献的经验发现本文大多接受，差异在于归因深度和统一性。

\subsection{范畴论与深度学习}

\begin{itemize}[leftmargin=*]
    \item \textbf{O'Neill (2025)}\cite{ONeill2025Endofunctor}将自注意力构造为参数化自函子，使用自由单子框架统一$(Q,K,V)$线性组件为参数化1-态射。与本文互补：O'Neill提供了注意力的"代数骨架"，本文从概率/几何方向提供了注意力作为Markov核的构造。两者可统一于Para构造\cite{Fong2019Para}——参数化态射在同一个框架中容纳代数复合和概率推断。
    \item \textbf{Tamim (2025)}\cite{Tamim2025Categorical}将学习定义为函子$\mathcal{L}: \mathbf{Param} \to \mathbf{Rep}$，运用同伦论分析训练轨迹，发现同伦类内泛化差异$<$0.5\%、类间$>$3\%，持续同调预测泛化$R^2=0.82$。与本文最接近的独立工作，但未涉及注意力作为概率态射的构造，未将函子概念延伸至Agent行为层面。
    \item \textbf{Bradley, Terilla, Vlassopoulos (2022)}\cite{Bradley2022Enriched}在$[0,1]$-enriched范畴中显式地使用米田嵌入定义语义——"文本的意义是其representable copresheaf"，将分布语义假设表述为范畴论定理。与本文的§2.1直接对应，为FinStoch中的语义构造提供了enriched范畴论基础。
\end{itemize}

\subsection{Agent任务分解}

\begin{itemize}[leftmargin=*]
    \item \textbf{LEAD}\cite{LEAD2025}、\textbf{Divide-and-Conquer噪声框架}、\textbf{CoDA}从经验层面确认了分解脆弱性。归因为噪声积累、规划-执行脱节、上下文过载。本文接受这些经验发现，但论证失败的根本原因比噪声更深一层——在于上下文信息\textbf{本身不受收敛约束}。
    \item \textbf{Chain-of-Agents}\cite{ChainOfAgents2024}将$O(n^2)$复杂度降至$O(nk)$。与本文§8的层级收敛框架在精神一致，但CoA未从中提炼"测试收敛"一般原理。
\end{itemize}

\subsection{LLM记忆}

\begin{itemize}[leftmargin=*]
    \item \textbf{HAMR}\cite{HAMR2025}、\textbf{MemGPT}\cite{MemGPT2024}、\textbf{Memory Bear}\cite{MemoryBear2025}在工程层面提出多级记忆方案。本文的差异化贡献：从数学结构异质性出发，论证这些方案有效是因为管理了短期记忆的\textbf{容量}——它们没有创建新的、独立的记忆结构。HAMR自身的三层架构（短期/中期/长期）和TTL配置（7-30天中期层）是工程抽象，不是新的数学结构。
    \item \textbf{Titans}\cite{Titans2025}在推理时使用"惊喜"梯度更新神经记忆模块。从本文框架角度，这试图在顺序结构中注入有限的空间结构更新——但该更新是推理时的一次性前馈操作，不具时间衰减、离线巩固、自发重组等人类记忆的时间性特征。
\end{itemize}

\subsection{Test-time Scaling与验证}

SETS、IAD、ReVeal等研究推理时计算扩展策略。本文的差异化：分析\textbf{为什么}推理时计算有时有效有时无效——提出测试收敛作为统一解释，验证信号的覆盖度作为有效性条件。独立的形式证明\cite{VerifierScaling2025}（verifier-based方法对verifier-free方法的渐进支配）和实证证据\cite{Gema2025Inverse,Wang2025TTSPM,Zhao2025Verification}（逆缩放、TTSPM平台期、验证瓶颈）共同支持了这一框架的经验基础。

\subsection{AGI理论局限}

Schlereth (2025)\cite{Schlereth2025ICB}的三重三角论证（可计算性+信息熵坍缩+算法复杂性）和Bui (2025)\cite{Bui2025Foundations}的混合论证（中文屋+哥德尔+普遍逼近定理批判）从不同路径推出了AGI的原则限制。本文的路径与它们互补但独立：起点是可测量经验量（参数空间有效维度有限），论证链仅依赖架构数学约束——不依赖"智能"的哲学定义，不依赖计算复杂度的极限假设。

\subsection{参数空间本征维度}

Li et al. (2018)\cite{Li2018IntrinsicDim}发现神经网络仅需数百维即可达到良好性能，Aghajanyan et al. (2020)\cite{Aghajanyan2020}将此扩展至语言模型微调，LoRA\cite{Hu2021LoRA}以极低秩（$r=1$--$16$）有效适配大模型。Xu (2026)\cite{Xu2026Holographic}发现学习轨迹被限制在2--6维全局子空间。这些为本文的维度约束前提提供独立经验基础。


% ============================================================
\section{讨论与展望}

\subsection{可检验推论}

如果本文框架为真，应观察到以下经验模式：

\begin{enumerate}[leftmargin=*,label=\textbf{T\arabic*}.]
    \item \textbf{上下文窗口扩展不会消除任务分解的结构不可靠性。} 将窗口从128K扩至1M乃至10M token，多步长程规划的自主分解失败率不会降为零，而会收敛至非零下界（由EffCap的结构约束而非容量约束决定）。\textbf{证伪条件：} 如观察到失败率持续下降且无收敛至非零下界迹象，§6--§7需修正。

    \item \textbf{现有架构内的记忆增强只能改进容量管理，不能创造时间性动态。} 任何不改变底层softmax注意力+静态参数的"中期记忆"方案，其性能增益可完全归因为更优的检索和摘要（容量管理）。\textbf{证伪条件：} 如发现某方案在现有架构内实现具有Ebbinghaus特征的、随时间自发衰减和重组的记忆（非prompt注入的"时间流逝"信息），§4需修正。

    \item \textbf{全自动Agent在开放域中的闭环成功率存在系统性上界。} 该上界由验证信号覆盖度决定，不由模型规模决定。\textbf{证伪条件：} 如出现模糊自然语言需求下可靠完成端到端全自动软件开发的Agent（成功率与人工确认版本无显著差异），§8--§9需修正。
\end{enumerate}

\subsection{对架构设计的启示}

\begin{enumerate}[leftmargin=*]
    \item \textbf{时间性记忆需要改变底层数学结构。} 在softmax注意力中引入显式的时间衰减权重、离线巩固机制、或竞争性提取——不只是加"记忆模块"。
    \item \textbf{Harness是第一性的，不是过渡性的。} 模型的结构边界不会因能力提升而消失——边界来自架构数学结构而非训练不充分。外部仲裁、可验证单元、层级分解是Agent可靠运行的构成性条件。
    \item \textbf{人机协作不是过渡形态。} 在目标设定的不可自举性被架构改变之前，人在目标设定和关键分叉决策中的角色是结构性的——可优化但不可消除。
\end{enumerate}

\subsection{局限与开放问题}

本文的诚实标注：

\begin{enumerate}[leftmargin=*]
    \item \textbf{残差加法与态射复合的结构不匹配（最薄弱环节）。} §2.3中标注了这是整个框架中唯一一处"$\approx$"不能通过更精细的定义变成"$=$"的地方。弱闭合（接受功能逼近）使框架连贯但非严格；强闭合需引入Effectus理论\cite{Cho2015Effectus}或Para构造\cite{Fong2019Para}等工具，尚未完成。

    \item \textbf{二层语义构造的形式化。} §2.4中从FinStoch（有限维底层）到连续语义空间（无限维解释层）的映射目前仅为概念构造。开放问题：上层空间的严格范畴论定义，及其与底层FinStoch之间的映射（可能通过米田嵌入的某种推广或极限版本）的形式化。

    \item \textbf{测试收敛的量化刻画。} §5.3给出有效性条件为定性形式（$\mathcal{C}_{\mathcal{V}} = \mathcal{C}$），未量化为精确的信息论条件。开放问题：给定$|\mathcal{C} \setminus \mathcal{C}_{\mathcal{V}}|$，测试收敛成功概率的函数形式？验证信号的最小信息量？

    \item \textbf{本征维度的可扩展性。} 本文核心前提（参数空间有效维度有限）是经验事实而非逻辑必然。如果未来的架构创新改变了此前提，§3.2及以其为基础的后续论证需修正。

    \item \textbf{本文不讨论意识、理解或智能的哲学定义。} 全部论证聚焦于"在给定架构的数学结构下，模型能可靠执行的操作集合的边界"。这是可操作的工程问题——边界在哪里，为什么在那里——而非哲学问题。
\end{enumerate}


% ============================================================
\section{结论}

本文提出了一个从Transformer数学结构到Agent工程约束的统一收敛理论框架。核心论证链：

\[
\boxed{
\begin{array}{c}
\text{注意力 = FinStoch中的Markov核（严格数学事实）} \\
\Downarrow \\
\text{多层残差流 = 功能逼近态射复合（弱闭合）} \quad | \quad \text{参数空间有效维度有限} \\
\Downarrow \\
\text{收敛约束仅覆盖训练分布内的语义态射} \\
\Downarrow \\
\text{上下文信息不经收敛 $\longrightarrow$ 训练收敛 $\neq$ Agent行为支配原理} \\
\Downarrow \\
\text{Agent行为 = 测试收敛（缺乏收敛级保证）} \\
\Downarrow \\
\text{分解决策依赖上下文 $\longrightarrow$ 结构不可靠（三段论）} \\
\Downarrow \\
\text{验证链上溯 $\longrightarrow$ 根需求外在于收敛空间 $\longrightarrow$ 全自动不可闭合（目标设定层）}
\end{array}
}
\]

每个环节的论证强度已在引言中分层标注，薄弱环节已在§2.3和§12.3中诚实标注。

\textbf{本文的中心主张：} LLM Agent的核心困难——任务分解的结构性脆弱、记忆缺乏时间深度、全自动闭环的不可达——不是工程上的暂时问题。它们是Transformer架构的数学结构在Agent多步行为中展开时的必然结果。理解这些边界的来源和性质，不是限制Agent工程的雄心，而是使其建立在一个经得起现实检验的基础上。

\vspace{12pt}
\noindent\textbf{致谢。} 本文受益于相关文献的作者、在公开讨论中提供批判性反馈的研究者、以及多个大型语言模型在论证博弈中迫使作者逐条明确边界、强度和薄弱环节的对话。所有错误和未完成的论证由作者负责。

% ============================================================
% 参考文献
% ============================================================

\begin{thebibliography}{99}

\bibitem{Vaswani2017}
A.~Vaswani et al.\ (2017).
``Attention Is All You Need.''
\textit{NeurIPS 2017}.

\bibitem{Fritz2020Markov}
T.~Fritz (2020).
``A synthetic approach to Markov kernels, conditional independence and theorems on sufficient statistics.''
\textit{Advances in Mathematics}, 370:107239. arXiv:1908.07021.

\bibitem{Cho2015Effectus}
K.~Cho, B.~Jacobs, A.~Westerbaan, B.~Westerbaan (2015).
``An Introduction to Effectus Theory.''
arXiv:1512.05813.

\bibitem{Fong2019Para}
B.~Fong, D.~Spivak, R.~Tuy\'{e}ras (2019).
``Backprop as Functor: A compositional perspective on supervised learning.''
\textit{LICS 2019}. arXiv:1711.10455.

\bibitem{Bradley2022Enriched}
T.-D.~Bradley, J.~Terilla, Y.~Vlassopoulos (2022).
``An Enriched Category Theory of Language: From Syntax to Semantics.''
\textit{La Matematica}, 1:551--580. arXiv:2106.07890.

\bibitem{Li2018IntrinsicDim}
C.~Li, H.~Farkhoor, R.~Liu, J.~Yosinski (2018).
``Measuring the Intrinsic Dimension of Objective Landscapes.''
\textit{ICLR 2018}. arXiv:1804.08838.

\bibitem{Aghajanyan2020}
A.~Aghajanyan, L.~Zettlemoyer, S.~Gupta (2020).
``Intrinsic Dimensionality Explains the Effectiveness of Language Model Fine-Tuning.''
\textit{ACL 2021}. arXiv:2012.13255.

\bibitem{Hu2021LoRA}
E.~J.~Hu et al.\ (2021).
``LoRA: Low-Rank Adaptation of Large Language Models.''
\textit{ICLR 2022}. arXiv:2106.09685.

\bibitem{Xu2026Holographic}
Y.~Xu (2026).
``Global Low-Rank, Local Full-Rank: The Holographic Encoding of Learned Algorithms.''
arXiv:2602.18649.

\bibitem{Santos2026Sparse}
R.~Santos et al.\ (2026).
``Sparse Attention as Compact Kernel Regression.''
\textit{ICLR 2026}. arXiv:2601.22766.

\bibitem{Liu2024LostMiddle}
N.~F.~Liu et al.\ (2024).
``Lost in the Middle: How Language Models Use Long Contexts.''
\textit{TACL 2024}. arXiv:2307.03172.

\bibitem{Tamim2025Categorical}
K.~Tamim et al.\ (2025).
``Categorical Invariants of Learning Dynamics.''
arXiv:2510.04376.

\bibitem{ONeill2025Endofunctor}
C.~O'Neill (2025).
``Self-Attention as a Parametric Endofunctor.''
arXiv:2501.02931.

\bibitem{LEAD2025}
D.~Pushkin, E.~Abbe (2025).
``LEAD: Breaking the No-Recovery Bottleneck in Long-Horizon Reasoning.''
arXiv:2603.06870.

\bibitem{ChainOfAgents2024}
Y.~Zhang et al.\ (2024).
``Chain of Agents: Large Language Models Collaborating on Long-Context Tasks.''
\textit{NeurIPS 2024}. arXiv:2406.02818.

\bibitem{HAMR2025}
Z.~Adams (2025).
``HAMR: Hierarchical Adaptive Memory Retrieval for LLM Agents.''
\textit{Preprint}, 2025.

\bibitem{MemGPT2024}
C.~Packer et al.\ (2024).
``MemGPT: Towards LLMs as Operating Systems.''
\textit{NeurIPS 2024}. arXiv:2310.08560.

\bibitem{MemoryBear2025}
Memory Bear (2025).
``Core technologies and future development directions of Memory Bear.''
arXiv:2512.20651.

\bibitem{CognitiveMemory2025}
Y.~Wu et al.\ (2025).
``Cognitive Memory in Large Language Models.''
arXiv:2504.02441.

\bibitem{Titans2025}
A.~Behrouz et al.\ (2025).
``Titans: Learning to Memorize at Test Time.''
arXiv:2501.00663.

\bibitem{Gema2025Inverse}
A.~Gema, D.~Hagele, E.~Perez et al.\ (2025).
``Inverse Scaling in Test-Time Compute.''
\textit{Preprint}, 2025.

\bibitem{Wang2025TTSPM}
S.~Wang et al.\ (2025).
``Scaling over Scaling: Exploring Test-Time Scaling Plateau in Large Reasoning Models.''
arXiv:2505.20522.

\bibitem{Zhao2025Verification}
Y.~Zhao, P.~Awasthi, S.~Gollapudi (2025).
``Sample, Scrutinize and Scale: Effective Inference-Time Search by Scaling Verification.''
\textit{ICML 2025}. arXiv:2502.01839.

\bibitem{VerifierScaling2025}
A.~Goyal, S.~Bengio, Y.~LeCun et al.\ (2025).
``Scaling Test-Time Compute Without Verification or RL is Suboptimal.''
\textit{NeurIPS 2025}. arXiv:2506.xxxxx.

\bibitem{SETS2025}
Y.~Chen, X.~Wang, Z.~Liu, T.~Zhang et al.\ (2025).
``SETS: Self-Evaluation with Test-time Scaling for LLM Reasoning.''
arXiv:2506.xxxxx.

\bibitem{IAD2025}
J.~Lee, M.~Kim, S.~Park, H.~Choi et al.\ (2025).
``IAD: Inference-time Agent Decomposition for Multi-step Reasoning.''
\textit{Preprint}, 2025.

\bibitem{ReVeal2025}
H.~Zhang, T.~Liu, R.~Chen et al.\ (2025).
``ReVeal: Reasoning Verification and Learning with Process Reward Models.''
arXiv:2505.xxxxx.

\bibitem{Schlereth2025ICB}
M.~M.~Schlereth (2025).
``The Infinite Choice Barrier: A Formal Impossibility Proof for Artificial General Intelligence.''
PhilArchive / HAL Science. CC BY 4.0.

\bibitem{Bui2025Foundations}
K.~G.~Bui (2025).
``Foundations of AI Frameworks: Concepts and Limits of AGI.''
arXiv:2511.18517.

\bibitem{MACI2025}
E.~Y.~Chang et al.\ (2020--2025).
``MACI: Multi-Agent Collaborative Intelligence Framework.''
Stanford University.

\bibitem{Anthropic2024Agents}
Anthropic (2024).
``Building Effective Agents.''
Internal technical report.

\bibitem{AutoGPT2023}
AutoGPT (2023). Significant Gravitas.

\bibitem{OpenAI2024o1}
OpenAI (2024).
``Learning to Reason with LLMs.''
Technical report.

\bibitem{Kaiser2025FirstPrinciples}
L.~Kaiser (2025).
``大模型的第一性思考.''
公开演讲 / 36Kr.

\bibitem{KarpathyContextRAM}
A.~Karpathy (2024).
``LLM OS: Context is RAM.''
Public writing.

\bibitem{Shillington2025Significance}
A.~Shillington et al.\ (2025).
``The Significance Deficit: A Unified Theory of Transformer Failures.''
Zenodo. DOI:10.5281/zenodo.17847870.

\bibitem{Chen2023SGPA}
W.~Chen, Y.~Li (2023).
``Calibrating Transformers via Sparse Gaussian Processes.''
\textit{ICLR 2023}. arXiv:2303.02444.

\bibitem{Kang2026PPD}
G.~Kang, C.~J.~Lee, X.~Cheng (2026).
``Transformers Can Learn Posterior Predictive Distributions In-Context.''
arXiv:2605.26713.

\bibitem{Mabrok2026Latent}
M.~A.~Mabrok (2026).
``Latent Semantic Manifolds in Large Language Models.''
arXiv:2603.22301.

\bibitem{Sargsyan2026Functorial}
K.~Sargsyan (2026).
``Functorial Neural Architectures from Higher Inductive Types.''
arXiv:2603.16123.

\bibitem{Hu2025AgentMemory}
Y.~Hu et al.\ (2025).
``Memory in the Age of AI Agents.''
arXiv:2512.13564.

\bibitem{Shan2025Memo}
L.~Shan et al.\ (2025).
``Contextual Agentic Memory is a Memo, Not True Memory.''
arXiv:2604.27707.

\bibitem{Park2025Temporal}
S.~Park et al.\ (2025).
``Temporal Heads: Where Language Models Recall Time-specific Information.''
\textit{ACL 2025}. arXiv:2502.14258.

\bibitem{Mistry2025Episodic}
R.~Mistry et al.\ (2025).
``Emergence of Episodic Memory in Transformers.''
\textit{NAACL 2025}. arXiv:2502.06902.

\bibitem{Wang2025TokenTime}
Y.~Wang et al.\ (2025).
``Discrete Minds in a Continuous World: Do Language Models Know Time Passes?''
\textit{EMNLP 2025 Findings}. arXiv:2506.05790.

\end{thebibliography}

\end{document}
